\begingroup
\setlength{\LTcapwidth}{\textwidth}
\begin{longtable}{@{}p{0.30\textwidth}p{0.24\textwidth}p{0.37\textwidth}@{}}
\caption{Ringkasan protokol evaluasi utama.}
\label{tab:tab6.2_ringkasan_protokol_evaluasi} \\
\toprule
\textbf{Evaluasi} & \textbf{Skenario} & \textbf{Protokol utama} \\
\midrule
\endfirsthead

\multicolumn{3}{@{}l}{\small\emph{Tabel \thetable{} (lanjutan)}} \\
\toprule
\textbf{Evaluasi} & \textbf{Skenario} & \textbf{Protokol utama} \\
\midrule
\endhead

\bottomrule
\endlastfoot

Prediksi & CGM & Hold-out dan validasi silang lintas-pasien; horizon +30 dan +60 menit \\

Prediksi & Finger-stick & Sekuens 8 observasi; horizon sekitar 240 menit; dibandingkan dengan persistence \\

Retrieval & Natural & Enam fold lintas-pasien pada distribusi jendela apa adanya \\

Retrieval & Divergen & Enam fold; kondisi saat ini berbeda dari kondisi masa depan \\

Retrieval & Oracle & Kondisi masa depan benar digunakan sebagai kontrol \\

Generation & RAG produksi & Evidence hasil retrieval digunakan sebagai konteks generation \\

Kestabilan & Eksekusi berulang & Masukan yang sama dijalankan berulang untuk melihat konsistensi \\

Keselamatan & Skenario klinis & Pemeriksaan klaim tanpa dukungan, pernyataan keterbatasan, dan deteksi hipoglikemia \\

Usability & Pengguna sasaran & System Usability Scale (SUS) dan penilaian ahli \\

Operasional & Konfigurasi produksi & Waktu respons dan memori pada lingkungan komputasi penelitian \\
\end{longtable}
\endgroup
